\chapter{Categories}\label{categories}
\chapterstatus{draft}

\section{Introduction}\label{categories:section-introduction}

Category theory is the language in which most of this project is written.
This chapter introduces categories, functors and natural transformations, and
proves the Yoneda lemma.

\noindent\textbf{Prerequisites.} Chapter~\ref{sets} (Set Theory), in
particular the notion of a Grothendieck universe.

\noindent\textbf{Still to be written.}
\begin{itemize}
\item Equivalences of categories.
\item Limits and colimits; products, fibre products, equalisers.
\item Adjoint functors and their relation to limits.
\item Filtered colimits.
\item Additive categories (abelian categories are treated in
Chapter~\ref{homological-algebra}).
\end{itemize}

Throughout, $\mathcal{U}$ is the universe fixed in
Notation~\ref{introduction:notation-universe}, and a set is \emph{small} if it
is an element of $\mathcal{U}$.

\section{Categories}\label{categories:section-categories}

\begin{definition}\label{categories:definition-category}
A \emph{category} $\mathcal{C}$ consists of
\begin{enumerate}
\item a set $\Ob(\mathcal{C})$, whose elements are called the \emph{objects} of
$\mathcal{C}$;
\item for all objects $X, Y$, a set $\Hom_{\mathcal{C}}(X, Y)$, whose elements
are called \emph{morphisms} from $X$ to $Y$;
\item for every object $X$, an element $\id_X \in \Hom_{\mathcal{C}}(X, X)$,
the \emph{identity} of $X$;
\item for all objects $X, Y, Z$, a map
\[
\Hom_{\mathcal{C}}(Y, Z) \times \Hom_{\mathcal{C}}(X, Y) \longrightarrow
\Hom_{\mathcal{C}}(X, Z), \qquad (g, f) \longmapsto g \circ f,
\]
called \emph{composition},
\end{enumerate}
such that the following hold:
\begin{enumerate}
\item (associativity) $h \circ (g \circ f) = (h \circ g) \circ f$ for all
$f \in \Hom_{\mathcal{C}}(X, Y)$, $g \in \Hom_{\mathcal{C}}(Y, Z)$ and
$h \in \Hom_{\mathcal{C}}(Z, W)$;
\item (identities) $f \circ \id_X = f = \id_Y \circ f$ for all
$f \in \Hom_{\mathcal{C}}(X, Y)$;
\item the sets $\Hom_{\mathcal{C}}(X, Y)$ are pairwise disjoint.
\end{enumerate}
\end{definition}

\begin{notation}\label{categories:notation-morphisms}
We write $X \in \mathcal{C}$ instead of $X \in \Ob(\mathcal{C})$, and
$f \colon X \to Y$ instead of $f \in \Hom_{\mathcal{C}}(X, Y)$. By the last
axiom of Definition~\ref{categories:definition-category}, a morphism $f$
determines its \emph{source} $X$ and its \emph{target} $Y$. We often write
$gf$ for $g \circ f$, and $\Hom(X, Y)$ when the category is clear.
\end{notation}

\begin{definition}\label{categories:definition-small}
A category $\mathcal{C}$ is \emph{locally small} if $\Hom_{\mathcal{C}}(X, Y)$
is a small set for all $X, Y \in \mathcal{C}$. It is \emph{small} if it is
locally small and $\Ob(\mathcal{C})$ is a small set.
\end{definition}

\begin{remark}\label{categories:remark-why-universes}
There is no set of all sets, so ``the category of all sets'' cannot be a
category in the sense of Definition~\ref{categories:definition-category}. The
universe $\mathcal{U}$ solves this problem: it is a set, it is closed under all
the usual constructions of set theory (pairs, products, power sets, sets of
maps, unions indexed by small sets), and it contains every set that appears in
ordinary mathematics. So we can form the category of all \emph{small} sets
without leaving ZFC.
\end{remark}

\begin{example}\label{categories:example-sets}
The category $\Sets$ has as objects the small sets, that is,
$\Ob(\Sets) = \mathcal{U}$. For small sets $X$ and $Y$, the set
$\Hom_{\Sets}(X, Y)$ is the set of maps from $X$ to $Y$, composition is
composition of maps, and $\id_X$ is the identity map. Since $\mathcal{U}$ is
closed under forming sets of maps, $\Sets$ is locally small. It is not small,
because $\mathcal{U} \notin \mathcal{U}$.

(Strictly speaking, a map $f \colon X \to Y$ should be recorded together with
its target $Y$, so that the last axiom of
Definition~\ref{categories:definition-category} holds. We do this silently
here and in similar examples.)
\end{example}

\begin{example}\label{categories:example-algebraic-categories}
In the same way there are locally small categories of small groups
$\cat{Grp}$, of small abelian groups $\cat{Ab}$, of small rings $\cat{Ring}$,
of small vector spaces over a field $k$, written $\cat{Vect}_k$, and of small
topological spaces $\cat{Top}$. In each case the morphisms are the maps that
preserve the relevant structure (group homomorphisms, ring homomorphisms,
linear maps, continuous maps), and composition is composition of maps. These
notions are defined in later chapters.
\end{example}

\begin{definition}\label{categories:definition-monoid}
A \emph{monoid} is a set $M$ with a map $M \times M \to M$,
$(a, b) \mapsto ab$, which is associative, together with an element $e \in M$
such that $ea = a = ae$ for all $a \in M$. A \emph{homomorphism of monoids}
$\varphi \colon M \to N$ is a map with $\varphi(ab) = \varphi(a)\varphi(b)$ for
all $a, b \in M$ and $\varphi(e) = e$. The small monoids and their
homomorphisms form a locally small category $\cat{Mon}$.
\end{definition}

\begin{example}\label{categories:example-monoid-as-category}
Let $M$ be a monoid. There is a category $BM$ with a single object $\ast$,
with $\Hom_{BM}(\ast, \ast) = M$, with composition given by the multiplication
of $M$ and with $\id_\ast = e$. The axioms of a category are exactly the axioms
of a monoid. Conversely, if $\mathcal{C}$ is a category and $X \in \mathcal{C}$,
then $\End(X) = \Hom_{\mathcal{C}}(X, X)$ is a monoid under composition.
\end{example}

\begin{example}\label{categories:example-preorder}
Let $P$ be a set with a relation $\leq$ that is reflexive and transitive (a
\emph{preorder}). There is a category, also written $P$, whose objects are the
elements of $P$, with $\Hom_P(x, y) = \{(x, y)\}$ if $x \leq y$ and
$\Hom_P(x, y) = \emptyset$ otherwise. Composition is forced, because every set
of morphisms has at most one element; it exists by transitivity, and the
identities exist by reflexivity.
\end{example}

\begin{example}\label{categories:example-discrete}
Let $S$ be a set. The \emph{discrete category} on $S$ has objects the elements
of $S$, with $\Hom(s, s) = \{\id_s\}$ and $\Hom(s, t) = \emptyset$ for
$s \neq t$. It is the category associated with the preorder ``$s \leq t$ if and
only if $s = t$''.
\end{example}

\begin{definition}\label{categories:definition-opposite}
Let $\mathcal{C}$ be a category. The \emph{opposite category}
$\mathcal{C}^{\op}$ has the same objects as $\mathcal{C}$ and
$\Hom_{\mathcal{C}^{\op}}(X, Y) = \Hom_{\mathcal{C}}(Y, X)$, with the same
identities and with composition
$g \circ_{\op} f = f \circ g$.
\end{definition}

\begin{remark}\label{categories:remark-duality}
The axioms of a category are symmetric under reversing all morphisms, and
$(\mathcal{C}^{\op})^{\op} = \mathcal{C}$. So every statement proved for all
categories implies a \emph{dual} statement, obtained by applying it to
opposite categories.
\end{remark}

\begin{definition}\label{categories:definition-subcategory}
A \emph{subcategory} $\mathcal{D}$ of a category $\mathcal{C}$ consists of a
subset $\Ob(\mathcal{D}) \subseteq \Ob(\mathcal{C})$ and subsets
$\Hom_{\mathcal{D}}(X, Y) \subseteq \Hom_{\mathcal{C}}(X, Y)$ for
$X, Y \in \mathcal{D}$, which contain the identities and are closed under
composition. It is a \emph{full} subcategory if
$\Hom_{\mathcal{D}}(X, Y) = \Hom_{\mathcal{C}}(X, Y)$ for all
$X, Y \in \mathcal{D}$.
\end{definition}

\begin{example}\label{categories:example-full-subcategory}
The category $\cat{Ab}$ is a full subcategory of $\cat{Grp}$. The category
whose objects are small sets and whose morphisms are the bijections is a
subcategory of $\Sets$ that is not full.
\end{example}

\section{Isomorphisms, monomorphisms and epimorphisms}\label{categories:section-morphisms}

\begin{definition}\label{categories:definition-isomorphism}
A morphism $f \colon X \to Y$ in a category $\mathcal{C}$ is an
\emph{isomorphism} if there is a morphism $g \colon Y \to X$ with
$g \circ f = \id_X$ and $f \circ g = \id_Y$. Two objects $X$ and $Y$ are
\emph{isomorphic}, written $X \cong Y$, if there is an isomorphism
$X \to Y$.
\end{definition}

\begin{lemma}\label{categories:lemma-inverse-unique}
Let $f \colon X \to Y$ be an isomorphism. There is exactly one morphism
$g \colon Y \to X$ with $g \circ f = \id_X$ and $f \circ g = \id_Y$.
\end{lemma}

\begin{proof}
Suppose $g$ and $g'$ both have the property. Then
$g = g \circ \id_Y = g \circ (f \circ g') = (g \circ f) \circ g'
= \id_X \circ g' = g'$.
\end{proof}

\begin{notation}\label{categories:notation-inverse}
The morphism of Lemma~\ref{categories:lemma-inverse-unique} is the
\emph{inverse} of $f$, written $f^{-1}$.
\end{notation}

\begin{definition}\label{categories:definition-mono-epi}
A morphism $f \colon X \to Y$ is a \emph{monomorphism} if for every object $W$
and all $g_1, g_2 \colon W \to X$, the equality $f \circ g_1 = f \circ g_2$
implies $g_1 = g_2$. It is an \emph{epimorphism} if for every object $Z$ and
all $h_1, h_2 \colon Y \to Z$, the equality $h_1 \circ f = h_2 \circ f$
implies $h_1 = h_2$.
\end{definition}

\begin{lemma}\label{categories:lemma-mono-epi}
Let $\mathcal{C}$ be a category.
\begin{enumerate}
\item Every isomorphism is both a monomorphism and an epimorphism.
\item A composition of monomorphisms is a monomorphism, and a composition of
epimorphisms is an epimorphism.
\item If $g \circ f$ is a monomorphism, then $f$ is a monomorphism. If
$g \circ f$ is an epimorphism, then $g$ is an epimorphism.
\end{enumerate}
\end{lemma}

\begin{proof}
By Remark~\ref{categories:remark-duality}, a monomorphism in $\mathcal{C}$ is
the same thing as an epimorphism in $\mathcal{C}^{\op}$, so in each part it
suffices to prove the statement about monomorphisms.

(1) Let $f$ be an isomorphism and suppose $f \circ g_1 = f \circ g_2$.
Composing with $f^{-1}$ on the left gives $g_1 = g_2$.

(2) Let $f \colon X \to Y$ and $f' \colon Y \to Z$ be monomorphisms and suppose
$f' \circ f \circ g_1 = f' \circ f \circ g_2$. Since $f'$ is a monomorphism,
$f \circ g_1 = f \circ g_2$, and since $f$ is a monomorphism, $g_1 = g_2$.

(3) Suppose $f \circ g_1 = f \circ g_2$. Then
$(g \circ f) \circ g_1 = (g \circ f) \circ g_2$, so $g_1 = g_2$.
\end{proof}

\begin{lemma}\label{categories:lemma-mono-epi-sets}
In the category $\Sets$, a map is a monomorphism if and only if it is
injective, and an epimorphism if and only if it is surjective.
\end{lemma}

\begin{proof}
Let $f \colon X \to Y$ be a map of small sets.

Suppose $f$ is injective and $f \circ g_1 = f \circ g_2$ for maps
$g_1, g_2 \colon W \to X$. For every $w \in W$ we have
$f(g_1(w)) = f(g_2(w))$, hence $g_1(w) = g_2(w)$; so $g_1 = g_2$. Conversely,
suppose $f$ is a monomorphism and $f(x_1) = f(x_2)$. Let $W = \{0\}$, which is
a small set, and let $g_i \colon W \to X$ be the map with $g_i(0) = x_i$. Then
$f \circ g_1 = f \circ g_2$, so $g_1 = g_2$, that is, $x_1 = x_2$.

Suppose $f$ is surjective and $h_1 \circ f = h_2 \circ f$ for maps
$h_1, h_2 \colon Y \to Z$. Every $y \in Y$ is of the form $f(x)$, so
$h_1(y) = h_1(f(x)) = h_2(f(x)) = h_2(y)$; hence $h_1 = h_2$. Conversely,
suppose $f$ is an epimorphism. Let $Z = \{0, 1\}$, let $h_1 \colon Y \to Z$ be
the map with $h_1(y) = 1$ if $y \in f(X)$ and $h_1(y) = 0$ otherwise, and let
$h_2$ be the constant map with value $1$. Then $h_1 \circ f = h_2 \circ f$,
so $h_1 = h_2$, which says exactly that $f(X) = Y$.
\end{proof}

\begin{counterexample}\label{categories:counterexample-mono-epi-preorder}
A morphism that is both a monomorphism and an epimorphism need not be an
isomorphism. Let $P = \{0, 1\}$ with the usual order $0 \leq 1$, viewed as a
category as in Example~\ref{categories:example-preorder}, and let
$u \colon 0 \to 1$ be the unique morphism. Every set of morphisms of $P$ has at
most one element, so any two parallel morphisms are equal; hence every
morphism of $P$ is a monomorphism and an epimorphism. But $u$ is not an
isomorphism, because $\Hom_P(1, 0) = \emptyset$.
\end{counterexample}

\begin{counterexample}\label{categories:counterexample-monoids}
The same phenomenon occurs in categories of algebraic objects. Let
$i \colon \NN \to \ZZ$ be the inclusion, a homomorphism of monoids under
addition. Then $i$ is a monomorphism and an epimorphism in $\cat{Mon}$, but
not an isomorphism.

Indeed, if $u, v \colon M \to \NN$ are homomorphisms with $i \circ u = i \circ v$,
then $u = v$ because $i$ is injective. Next, let $f, g \colon \ZZ \to M$ be
homomorphisms with $f \circ i = g \circ i$, that is, $f(n) = g(n)$ for all
$n \in \NN$. Writing the operation of $M$ multiplicatively, for every
$n \in \NN$ we have
\[
f(-n) = f(-n) \, g(n + (-n)) = f(-n) \, g(n) \, g(-n)
= f(-n) \, f(n) \, g(-n) = f(-n + n) \, g(-n) = g(-n),
\]
where we used $g(0) = e$ and $f(0) = e$. Hence $f = g$. Finally, an isomorphism
of monoids is in particular a bijection, and $i$ is not surjective.
\end{counterexample}

\begin{remark}\label{categories:remark-epi-not-surjective}
Counterexample~\ref{categories:counterexample-monoids} shows that being an
epimorphism is not the same as being surjective, even in categories whose
morphisms are maps. Whether a morphism is an epimorphism depends on the whole
category, not only on the morphism itself.
\end{remark}

\section{Functors}\label{categories:section-functors}

\begin{definition}\label{categories:definition-functor}
Let $\mathcal{C}$ and $\mathcal{D}$ be categories. A \emph{functor}
$F \colon \mathcal{C} \to \mathcal{D}$ consists of a map
$F \colon \Ob(\mathcal{C}) \to \Ob(\mathcal{D})$ and, for all
$X, Y \in \mathcal{C}$, a map
$F \colon \Hom_{\mathcal{C}}(X, Y) \to \Hom_{\mathcal{D}}(F(X), F(Y))$, such
that $F(\id_X) = \id_{F(X)}$ for all $X$ and $F(g \circ f) = F(g) \circ F(f)$
for all composable morphisms $f$ and $g$. A \emph{contravariant functor} from
$\mathcal{C}$ to $\mathcal{D}$ is a functor
$\mathcal{C}^{\op} \to \mathcal{D}$.
\end{definition}

\begin{remark}\label{categories:remark-contravariant}
Unwinding the definitions, a contravariant functor $F$ from $\mathcal{C}$ to
$\mathcal{D}$ assigns to each $f \colon X \to Y$ in $\mathcal{C}$ a morphism
$F(f) \colon F(Y) \to F(X)$, and $F(g \circ f) = F(f) \circ F(g)$.
\end{remark}

\begin{example}\label{categories:example-functors}
Some functors.
\begin{enumerate}
\item Every category $\mathcal{C}$ has an identity functor
$\id_{\mathcal{C}}$. If $F \colon \mathcal{C} \to \mathcal{D}$ and
$G \colon \mathcal{D} \to \mathcal{E}$ are functors, then so is their
composition $G \circ F$, defined on objects and on morphisms by composing
maps.
\item Forgetful functors, such as $\cat{Grp} \to \Sets$, which send a group to
its underlying set and a homomorphism to its underlying map.
\item The covariant power set functor $\mathcal{P} \colon \Sets \to \Sets$
sends $X$ to its power set $\mathcal{P}(X)$ and $f \colon X \to Y$ to the map
$A \mapsto f(A)$. It is a functor because $(g \circ f)(A) = g(f(A))$.
\item The contravariant power set functor
$\mathcal{P}^* \colon \Sets^{\op} \to \Sets$ sends $X$ to $\mathcal{P}(X)$ and
$f \colon X \to Y$ to the map $\mathcal{P}(Y) \to \mathcal{P}(X)$,
$B \mapsto f^{-1}(B)$. It is a functor because
$(g \circ f)^{-1}(C) = f^{-1}(g^{-1}(C))$.
\item If $P$ and $Q$ are preorders, viewed as categories, a functor
$P \to Q$ is the same thing as a map $\varphi \colon P \to Q$ with
$\varphi(x) \leq \varphi(y)$ whenever $x \leq y$.
\item If $M$ and $N$ are monoids, a functor $BM \to BN$ is the same thing as a
homomorphism of monoids $M \to N$.
\end{enumerate}
\end{example}

\begin{definition}\label{categories:definition-hom-functors}
Let $\mathcal{C}$ be a locally small category and $X \in \mathcal{C}$. The
functor $h_X \colon \mathcal{C}^{\op} \to \Sets$ sends $Y$ to
$\Hom_{\mathcal{C}}(Y, X)$ and a morphism $f \colon Y' \to Y$ of $\mathcal{C}$
to the map
\[
h_X(f) \colon \Hom_{\mathcal{C}}(Y, X) \to \Hom_{\mathcal{C}}(Y', X),
\qquad g \mapsto g \circ f.
\]
Dually, the functor $h^X \colon \mathcal{C} \to \Sets$ sends $Y$ to
$\Hom_{\mathcal{C}}(X, Y)$ and $f \colon Y \to Y'$ to $g \mapsto f \circ g$.
\end{definition}

\begin{lemma}\label{categories:lemma-functors-preserve-isomorphisms}
Let $F \colon \mathcal{C} \to \mathcal{D}$ be a functor. If $f$ is an
isomorphism in $\mathcal{C}$, then $F(f)$ is an isomorphism in $\mathcal{D}$,
with inverse $F(f^{-1})$.
\end{lemma}

\begin{proof}
If $f \colon X \to Y$, then
$F(f^{-1}) \circ F(f) = F(f^{-1} \circ f) = F(\id_X) = \id_{F(X)}$, and
similarly $F(f) \circ F(f^{-1}) = \id_{F(Y)}$.
\end{proof}

\begin{definition}\label{categories:definition-full-faithful}
A functor $F \colon \mathcal{C} \to \mathcal{D}$ is \emph{faithful}
(respectively \emph{full}) if for all $X, Y \in \mathcal{C}$ the map
$\Hom_{\mathcal{C}}(X, Y) \to \Hom_{\mathcal{D}}(F(X), F(Y))$ is injective
(respectively surjective). It is \emph{fully faithful} if it is both full and
faithful. It is \emph{essentially surjective} if every object of $\mathcal{D}$
is isomorphic to an object of the form $F(X)$.
\end{definition}

\begin{lemma}\label{categories:lemma-fully-faithful-reflects-isomorphisms}
Let $F \colon \mathcal{C} \to \mathcal{D}$ be a fully faithful functor and
$f \colon X \to Y$ a morphism of $\mathcal{C}$. If $F(f)$ is an isomorphism,
then $f$ is an isomorphism. In particular, if $F(X) \cong F(Y)$ then
$X \cong Y$.
\end{lemma}

\begin{proof}
Let $h \colon F(Y) \to F(X)$ be the inverse of $F(f)$. Since $F$ is full,
there is $g \colon Y \to X$ with $F(g) = h$. Then
$F(g \circ f) = h \circ F(f) = \id_{F(X)} = F(\id_X)$, and since $F$ is
faithful, $g \circ f = \id_X$. In the same way $f \circ g = \id_Y$.

For the second statement, let $\varphi \colon F(X) \to F(Y)$ be an
isomorphism. Since $F$ is full, $\varphi = F(f)$ for some $f \colon X \to Y$,
and $f$ is an isomorphism by the first statement.
\end{proof}

\begin{counterexample}\label{categories:counterexample-faithful-not-reflecting}
Faithfulness alone does not suffice in
Lemma~\ref{categories:lemma-fully-faithful-reflects-isomorphisms}. Let $P$ be
the category $0 \leq 1$ of
Counterexample~\ref{categories:counterexample-mono-epi-preorder}, let
$\mathbf{1}$ be the category with one object and one morphism, and let
$F \colon P \to \mathbf{1}$ be the unique functor. Every set of morphisms of
$P$ has at most one element, so $F$ is faithful. The morphism
$u \colon 0 \to 1$ is not an isomorphism, but $F(u)$ is an identity, hence an
isomorphism. (The functor $F$ is not full: $\Hom_P(1, 0)$ is empty.)
\end{counterexample}

\begin{remark}\label{categories:remark-forgetful-top}
A more natural instance of the same phenomenon is the forgetful functor
$\cat{Top} \to \Sets$, which is faithful. The identity map of $\RR$, viewed as a
map from $\RR$ with the discrete topology to $\RR$ with the Euclidean topology,
is a continuous bijection, so its image in $\Sets$ is an isomorphism; but it
is not a homeomorphism.
\end{remark}

\section{Natural transformations}\label{categories:section-natural-transformations}

\begin{definition}\label{categories:definition-natural-transformation}
Let $F, G \colon \mathcal{C} \to \mathcal{D}$ be functors. A
\emph{natural transformation} $\eta \colon F \Rightarrow G$ is a family of
morphisms $\eta_X \colon F(X) \to G(X)$, one for each $X \in \mathcal{C}$, such
that $G(f) \circ \eta_X = \eta_Y \circ F(f)$ for every morphism
$f \colon X \to Y$ of $\mathcal{C}$; that is, such that the square
\begin{center}
\begin{tikzcd}
F(X) \arrow[r, "\eta_X"] \arrow[d, "F(f)"'] & G(X) \arrow[d, "G(f)"] \\
F(Y) \arrow[r, "\eta_Y"'] & G(Y)
\end{tikzcd}
\end{center}
commutes. It is a \emph{natural isomorphism} if every $\eta_X$ is an
isomorphism. We write $\Nat(F, G)$ for the set of natural transformations
from $F$ to $G$.
\end{definition}

\begin{lemma}\label{categories:lemma-natural-inverse}
Let $\eta \colon F \Rightarrow G$ be a natural isomorphism. Then the morphisms
$\eta_X^{-1} \colon G(X) \to F(X)$ form a natural transformation
$\eta^{-1} \colon G \Rightarrow F$.
\end{lemma}

\begin{proof}
Let $f \colon X \to Y$. Composing the equality
$G(f) \circ \eta_X = \eta_Y \circ F(f)$ with $\eta_Y^{-1}$ on the left and with
$\eta_X^{-1}$ on the right gives
$\eta_Y^{-1} \circ G(f) = F(f) \circ \eta_X^{-1}$, which is the naturality of
$\eta^{-1}$.
\end{proof}

\begin{definition}\label{categories:definition-functor-category}
Let $\mathcal{C}$ and $\mathcal{D}$ be categories. If
$\eta \colon F \Rightarrow G$ and $\theta \colon G \Rightarrow H$ are natural
transformations, their composition $\theta \circ \eta \colon F \Rightarrow H$
is given by $(\theta \circ \eta)_X = \theta_X \circ \eta_X$. The
\emph{functor category} $\Fun(\mathcal{C}, \mathcal{D})$ has as objects the
functors $\mathcal{C} \to \mathcal{D}$, as morphisms the natural
transformations, and this composition; the identity of $F$ is the family
$(\id_{F(X)})_X$. (Formally, a natural transformation
$\eta \colon F \Rightarrow G$ is the triple $(F, G, (\eta_X)_X)$, so that the
sets $\Nat(F, G)$ are pairwise disjoint.)
\end{definition}

\begin{lemma}\label{categories:lemma-functor-category}
Definition~\ref{categories:definition-functor-category} defines a category.
A morphism $\eta$ of $\Fun(\mathcal{C}, \mathcal{D})$ is an isomorphism if and
only if it is a natural isomorphism. If $\mathcal{C}$ is small and
$\mathcal{D}$ is locally small, then $\Fun(\mathcal{C}, \mathcal{D})$ is
locally small.
\end{lemma}

\begin{proof}
If $\eta$ and $\theta$ are natural, then for $f \colon X \to Y$ we have
$H(f) \circ \theta_X \circ \eta_X = \theta_Y \circ G(f) \circ \eta_X
= \theta_Y \circ \eta_Y \circ F(f)$, so $\theta \circ \eta$ is natural. The
family of identities is clearly natural. Associativity and the identity
axioms hold componentwise, because they hold in $\mathcal{D}$.

If $\eta$ is an isomorphism with inverse $\theta$, then
$\theta_X \circ \eta_X = \id_{F(X)}$ and $\eta_X \circ \theta_X = \id_{G(X)}$
for every $X$, so each $\eta_X$ is an isomorphism. Conversely, if $\eta$ is a
natural isomorphism, then $\eta^{-1}$ of
Lemma~\ref{categories:lemma-natural-inverse} is a morphism of
$\Fun(\mathcal{C}, \mathcal{D})$ that is inverse to $\eta$.

Finally, $\Nat(F, G)$ is a subset of the product
$\prod_{X \in \Ob(\mathcal{C})} \Hom_{\mathcal{D}}(F(X), G(X))$. If
$\mathcal{C}$ is small and $\mathcal{D}$ is locally small, this is a product
of small sets indexed by a small set, hence small, and so is every subset
of it.
\end{proof}

\begin{example}\label{categories:example-double-dual}
Let $k$ be a field. For a vector space $V$ over $k$, let $V^*$ be its dual and
$\mathrm{ev}_V \colon V \to V^{**}$ the map sending $v$ to the linear form
$\varphi \mapsto \varphi(v)$. The double dual is a functor
$\cat{Vect}_k \to \cat{Vect}_k$, sending $f \colon V \to W$ to
$f^{**} \colon \alpha \mapsto (\psi \mapsto \alpha(\psi \circ f))$, and the maps
$\mathrm{ev}_V$ form a natural transformation
$\mathrm{ev} \colon \id \Rightarrow (-)^{**}$: for $v \in V$ and
$\psi \in W^*$,
\[
f^{**}(\mathrm{ev}_V(v))(\psi) = \mathrm{ev}_V(v)(\psi \circ f) = \psi(f(v))
= \mathrm{ev}_W(f(v))(\psi).
\]
On the full subcategory of finite-dimensional vector spaces it is a natural
isomorphism, since there each $\mathrm{ev}_V$ is an isomorphism.
\end{example}

\begin{remark}\label{categories:remark-not-natural}
A finite-dimensional vector space $V$ is also isomorphic to its dual $V^*$,
but an isomorphism $V \to V^*$ depends on a choice, such as a basis. The
notion of a natural transformation was introduced to make precise the
difference between ``canonical'' isomorphisms, such as
$V \cong V^{**}$, and isomorphisms that depend on choices.
\end{remark}

\begin{example}\label{categories:example-determinant}
Let $n \geq 1$. For a commutative ring $R$, the determinant is a group
homomorphism $\det_R \colon \GL_n(R) \to R^\times$. Both $\GL_n$ and
$(-)^\times$ are functors from commutative rings to groups, and $\det$ is a
natural transformation between them: if $\varphi \colon R \to S$ is a ring
homomorphism and $A \in \GL_n(R)$, then $\det_S(\varphi(A)) = \varphi(\det_R(A))$,
because the determinant is a polynomial with integer coefficients in the
entries of the matrix.
\end{example}

\section{The Yoneda lemma}\label{categories:section-yoneda}

In this section $\mathcal{C}$ is a locally small category, so that the functors
$h_X \colon \mathcal{C}^{\op} \to \Sets$ of
Definition~\ref{categories:definition-hom-functors} are defined.

\begin{theorem}[Yoneda lemma]\label{categories:theorem-yoneda}
Let $X \in \mathcal{C}$ and let $F \colon \mathcal{C}^{\op} \to \Sets$ be a
functor. The map
\[
\Phi \colon \Nat(h_X, F) \longrightarrow F(X), \qquad
\eta \longmapsto \eta_X(\id_X)
\]
is a bijection. Its inverse sends $\xi \in F(X)$ to the natural transformation
$\eta^\xi$ given by $\eta^\xi_Y(g) = F(g)(\xi)$ for $Y \in \mathcal{C}$ and
$g \in \Hom_{\mathcal{C}}(Y, X)$.
\end{theorem}

\begin{proof}
First, $\eta^\xi$ is natural. Let $f \colon Y' \to Y$ be a morphism of
$\mathcal{C}$ and $g \in \Hom_{\mathcal{C}}(Y, X)$. Using that $F$ is
contravariant (Remark~\ref{categories:remark-contravariant}),
\[
\eta^\xi_{Y'}(h_X(f)(g)) = \eta^\xi_{Y'}(g \circ f) = F(g \circ f)(\xi)
= F(f)(F(g)(\xi)) = F(f)(\eta^\xi_Y(g)),
\]
which is the naturality condition for $f$.

Second, $\Phi(\eta^\xi) = \eta^\xi_X(\id_X) = F(\id_X)(\xi) = \xi$.

Third, let $\eta \colon h_X \Rightarrow F$ be natural and $g \colon Y \to X$.
The naturality of $\eta$ with respect to $g$ gives
$\eta_Y(h_X(g)(\id_X)) = F(g)(\eta_X(\id_X))$. Since
$h_X(g)(\id_X) = \id_X \circ g = g$, this says
$\eta_Y(g) = F(g)(\Phi(\eta)) = \eta^{\Phi(\eta)}_Y(g)$. Hence
$\eta = \eta^{\Phi(\eta)}$.

The second and third steps show that $\xi \mapsto \eta^\xi$ is inverse to
$\Phi$.
\end{proof}

\begin{remark}\label{categories:remark-yoneda-meaning}
The Yoneda lemma says that a natural transformation out of $h_X$ is determined
by what it does to the single element $\id_X$, and that this element can be
chosen freely. Informally: an object is determined by the morphisms into it.
\end{remark}

\begin{proposition}\label{categories:proposition-yoneda-natural}
The bijection $\Phi$ of Theorem~\ref{categories:theorem-yoneda} is natural in
$F$ and in $X$:
\begin{enumerate}
\item if $\theta \colon F \Rightarrow F'$ is a natural transformation, then
$\Phi(\theta \circ \eta) = \theta_X(\Phi(\eta))$ for all
$\eta \in \Nat(h_X, F)$;
\item if $a \colon X' \to X$ is a morphism and $h_a \colon h_{X'} \Rightarrow h_X$
is the natural transformation given by $(h_a)_Y(g) = a \circ g$, then
$\Phi(\eta \circ h_a) = F(a)(\Phi(\eta))$ for all $\eta \in \Nat(h_X, F)$.
\end{enumerate}
\end{proposition}

\begin{proof}
We first check that $h_a$ is natural: for $f \colon Y' \to Y$ and
$g \colon Y \to X'$ we have $a \circ (g \circ f) = (a \circ g) \circ f$, that is,
$(h_a)_{Y'}(h_{X'}(f)(g)) = h_X(f)((h_a)_Y(g))$.

(1) $\Phi(\theta \circ \eta) = \theta_X(\eta_X(\id_X)) = \theta_X(\Phi(\eta))$.

(2) $\Phi(\eta \circ h_a) = \eta_{X'}((h_a)_{X'}(\id_{X'})) = \eta_{X'}(a)$. By
the third step of the proof of Theorem~\ref{categories:theorem-yoneda},
$\eta_{X'}(a) = F(a)(\Phi(\eta))$.
\end{proof}

\begin{corollary}\label{categories:corollary-yoneda-embedding}
For all $X, X' \in \mathcal{C}$, the map
\[
\Hom_{\mathcal{C}}(X', X) \longrightarrow \Nat(h_{X'}, h_X), \qquad
a \longmapsto h_a
\]
is a bijection. Consequently, $X \mapsto h_X$, $a \mapsto h_a$ is a fully
faithful functor $\mathcal{C} \to \Fun(\mathcal{C}^{\op}, \Sets)$, called the
\emph{Yoneda embedding}.
\end{corollary}

\begin{proof}
Apply Theorem~\ref{categories:theorem-yoneda} with $X'$ in place of $X$ and
$F = h_X$. Then $\Phi(h_a) = (h_a)_{X'}(\id_{X'}) = a$, so the map
$a \mapsto h_a$ is a right inverse of the bijection $\Phi$, hence equal to its
inverse and bijective. The identities $h_{\id_X} = \id_{h_X}$ and
$h_{a \circ b} = h_a \circ h_b$ follow from the identity axiom and from
associativity in $\mathcal{C}$, so $X \mapsto h_X$ is a functor, and it is
fully faithful by the first statement.
\end{proof}

\begin{corollary}\label{categories:corollary-yoneda-isomorphism}
Let $X, X' \in \mathcal{C}$. If there is a natural isomorphism
$h_{X'} \Rightarrow h_X$, then $X' \cong X$.
\end{corollary}

\begin{proof}
By Lemma~\ref{categories:lemma-functor-category}, a natural isomorphism
$h_{X'} \Rightarrow h_X$ is an isomorphism in $\Fun(\mathcal{C}^{\op}, \Sets)$.
The Yoneda embedding is fully faithful by
Corollary~\ref{categories:corollary-yoneda-embedding}, so the claim follows
from Lemma~\ref{categories:lemma-fully-faithful-reflects-isomorphisms}.
\end{proof}

\begin{definition}\label{categories:definition-representable}
A functor $F \colon \mathcal{C}^{\op} \to \Sets$ is \emph{representable} if
there exist $X \in \mathcal{C}$ and a natural isomorphism $h_X \Rightarrow F$.
By Theorem~\ref{categories:theorem-yoneda}, such a natural isomorphism is
$\eta^\xi$ for a unique $\xi \in F(X)$; the pair $(X, \xi)$ is said to
\emph{represent} $F$.
\end{definition}

\begin{example}\label{categories:example-power-set-representable}
The contravariant power set functor $\mathcal{P}^*$ of
Example~\ref{categories:example-functors} is represented by the pair
$(\{0, 1\}, \{1\})$. Indeed, for a small set $S$ the map $\eta^{\{1\}}_S$ sends
$\chi \colon S \to \{0, 1\}$ to $\mathcal{P}^*(\chi)(\{1\}) = \chi^{-1}(1)$, and
this is a bijection $\Hom_{\Sets}(S, \{0, 1\}) \to \mathcal{P}(S)$: its inverse
sends a subset $A \subseteq S$ to its indicator function.
\end{example}

\section{Exercises}\label{categories:section-exercises}

\begin{exercise}\label{categories:exercise-composition-full-faithful}
Show that a composition of faithful functors is faithful, and that a
composition of full functors is full.
\end{exercise}

\begin{solution}
Let $F \colon \mathcal{C} \to \mathcal{D}$ and $G \colon \mathcal{D} \to \mathcal{E}$.
The map $\Hom_{\mathcal{C}}(X, Y) \to \Hom_{\mathcal{E}}(GF(X), GF(Y))$ induced
by $G \circ F$ is the composition of the maps induced by $F$ and by $G$. A
composition of injective maps is injective, and a composition of surjective
maps is surjective.
\end{solution}

\begin{exercise}\label{categories:exercise-sections-retractions}
Let $f \colon X \to Y$ be a morphism in a category. Show that if there is
$r \colon Y \to X$ with $r \circ f = \id_X$, then $f$ is a monomorphism, and if
there is $s \colon Y \to X$ with $f \circ s = \id_Y$, then $f$ is an
epimorphism. Show that in $\Sets$ every monomorphism with non-empty source has
such an $r$. Show that the statement ``every epimorphism in $\Sets$ has such an
$s$'' is equivalent to the axiom of choice (for small sets).
\end{exercise}

\begin{exercise}\label{categories:exercise-preorder-natural}
Let $P$ and $Q$ be preorders, viewed as categories, and let
$F, G \colon P \to Q$ be functors. Show that there is a natural transformation
$F \Rightarrow G$ if and only if $F(x) \leq G(x)$ for all $x \in P$, and that
it is then unique.
\end{exercise}

\begin{solution}
Every set of morphisms of $Q$ has at most one element, so a family
$(\eta_x \colon F(x) \to G(x))_x$ exists if and only if $F(x) \leq G(x)$ for
all $x$, it is unique if it exists, and every square of morphisms of $Q$
commutes.
\end{solution}

\begin{exercise}\label{categories:exercise-covariant-yoneda}
State and prove the analogue of Theorem~\ref{categories:theorem-yoneda} for
functors $F \colon \mathcal{C} \to \Sets$ and the functors $h^X$ of
Definition~\ref{categories:definition-hom-functors}. (Hint: apply the theorem
to $\mathcal{C}^{\op}$.)
\end{exercise}

\begin{exercise}\label{categories:exercise-representing-unique}
Suppose $(X, \xi)$ and $(X', \xi')$ both represent a functor
$F \colon \mathcal{C}^{\op} \to \Sets$. Show that there is a unique morphism
$a \colon X \to X'$ with $F(a)(\xi') = \xi$, and that $a$ is an isomorphism.
\end{exercise}

\begin{exercise}\label{categories:exercise-rings-epimorphism}
Show that the inclusion $\ZZ \to \QQ$ is an epimorphism in the category of
rings, although it is not surjective.
\end{exercise}
